\documentclass[11pt,a4paper]{article}
\usepackage[utf8]{inputenc}
\usepackage[T1]{fontenc}
\usepackage{geometry}
\usepackage{hyperref}
\usepackage{listings}
\usepackage{xcolor}
\usepackage{booktabs}
\usepackage{enumitem}
\usepackage{fancyhdr}
\usepackage{titlesec}

% Page setup
\geometry{margin=1in}
\pagestyle{fancy}
\fancyhf{}
\fancyhead[L]{CEP Report: MLFQ Scheduler Implementation}
\fancyhead[R]{\thepage}
\fancyfoot[C]{}

% Code listing style
\lstset{
    language=C,
    basicstyle=\ttfamily\small,
    keywordstyle=\color{blue}\bfseries,
    commentstyle=\color{green!60!black},
    stringstyle=\color{red},
    numbers=left,
    numberstyle=\tiny\color{gray},
    stepnumber=1,
    numbersep=5pt,
    frame=single,
    breaklines=true,
    showstringspaces=false,
    tabsize=2
}

% Title formatting
\titleformat{\section}{\Large\bfseries}{\thesection}{1em}{}
\titleformat{\subsection}{\large\bfseries}{\thesubsection}{1em}{}
\titleformat{\subsubsection}{\normalsize\bfseries}{\thesubsubsection}{1em}{}

% Hyperref setup
\hypersetup{
    colorlinks=true,
    linkcolor=blue,
    filecolor=magenta,
    urlcolor=cyan,
    pdftitle={CEP Report: MLFQ Scheduler Implementation},
    pdfauthor={Your Name}
}

\begin{document}

% Title page
\begin{titlepage}
    \centering
    \vspace*{2cm}
    
    {\Huge\bfseries CEP Report}\\[0.5cm]
    {\Large Multi-Level Feedback Queue (MLFQ) Scheduler Implementation}\\[1cm]
    
    \vspace{1cm}
    
    {\large\textbf{Project}: Implementation of MLFQ Scheduler in xv6-riscv}\\[0.3cm]
    {\large\textbf{Date}: November 2024}\\[0.3cm]
    {\large\textbf{Status}: \textcolor{green!70!black}{Complete and Verified}}\\[1cm]
    
    \vfill
    
    % Optional: Uncomment and add your institution logo
    % \begin{center}
    %     \includegraphics[width=0.3\textwidth]{logo.png}
    % \end{center}
    
    \vspace{1cm}
    
    \today
\end{titlepage}

% Table of contents
\newpage
\tableofcontents
\newpage

% Introduction
\section{Introduction}

This project implements a Multi-Level Feedback Queue (MLFQ) scheduler for the xv6-riscv operating system. The MLFQ scheduler improves upon the default round-robin scheduler by providing priority-based scheduling that adapts to process behavior, giving higher priority to I/O-bound processes while preventing CPU-bound processes from starving lower-priority tasks.

\subsection{Objectives}

\begin{itemize}
    \item Replace the round-robin scheduler with an MLFQ scheduler
    \item Implement priority-based process scheduling with 4 priority levels
    \item Support automatic process demotion based on CPU usage
    \item Implement priority boost mechanism to prevent starvation
    \item Maintain system responsiveness for interactive processes
\end{itemize}

% Problem Statement
\section{Problem Statement}

The original xv6 scheduler uses a simple round-robin approach where all processes receive equal CPU time regardless of their characteristics. This leads to:

\begin{enumerate}
    \item \textbf{Poor responsiveness}: I/O-bound processes (like interactive applications) wait as long as CPU-bound processes
    \item \textbf{Inefficient scheduling}: No differentiation between process types
    \item \textbf{No priority management}: All processes treated equally
\end{enumerate}

The MLFQ scheduler addresses these issues by:
\begin{itemize}
    \item Assigning higher priority to I/O-bound processes
    \item Automatically demoting CPU-bound processes
    \item Providing periodic priority boosts to prevent starvation
\end{itemize}

% Design and Implementation
\section{Design and Implementation}

\subsection{MLFQ Architecture}

The implementation uses \textbf{4 priority queues} (numbered 0-3, where 0 is highest priority):

\begin{itemize}
    \item \textbf{Queue 0}: Highest priority, time allotment = 1 tick
    \item \textbf{Queue 1}: High priority, time allotment = 2 ticks
    \item \textbf{Queue 2}: Medium priority, time allotment = 4 ticks
    \item \textbf{Queue 3}: Lowest priority, time allotment = 8 ticks
\end{itemize}

\subsection{Key Components}

\subsubsection{Data Structures}

\textbf{Process Structure Extensions} (\texttt{kernel/proc.h}):
\begin{lstlisting}
// MLFQ Fields
int queue_level;             // Current Queue (0 = Highest Priority)
int time_in_queue;           // Ticks used in current queue
struct proc* next;           // Next Process in queue (For Linked List)
\end{lstlisting}

\textbf{Queue Structure} (\texttt{kernel/proc.c}):
\begin{lstlisting}
#define NMLFQ 4                 // Number of Queues
#define BOOST_INTERVAL 100      // Priority boost every 100 ticks

int time_allotment[NMLFQ] = {1, 2, 4, 8};

struct {
  struct proc* head;
  struct proc* tail;
  struct spinlock lock;
} mlfq[NMLFQ];
\end{lstlisting}

\subsubsection{Core Functions}

\begin{enumerate}
    \item \texttt{mlfq\_init()}: Initializes all queues and locks
    \item \texttt{mlfq\_enqueue(p, queue)}: Adds process to specified queue
    \item \texttt{mlfq\_dequeue(queue)}: Removes process from queue head
    \item \texttt{mlfq\_remove(p)}: Removes specific process from its queue
    \item \texttt{mlfq\_tick()}: Tracks time and handles priority boost
    \item \texttt{mlfq\_boost()}: Moves all processes to highest priority queue
\end{enumerate}

\subsection{Scheduling Algorithm}

\begin{enumerate}
    \item \textbf{Process Selection}: Scheduler checks queues from highest (0) to lowest (3) priority
    \item \textbf{Time Tracking}: Each timer interrupt increments \texttt{time\_in\_queue} for running process
    \item \textbf{Demotion}: When \texttt{time\_in\_queue >= time\_allotment[queue\_level]}, process is demoted
    \item \textbf{Priority Boost}: Every 100 ticks, all processes are boosted to queue 0
\end{enumerate}

\subsection{Process Lifecycle Integration}

\begin{itemize}
    \item \textbf{Fork/Userinit}: New processes start at queue 0 (highest priority)
    \item \textbf{Exit}: Processes removed from queue before termination
    \item \textbf{Sleep}: Processes removed from queue when sleeping
    \item \textbf{Wakeup}: Processes re-enqueued at current priority level
\end{itemize}

% Code Changes
\section{Code Changes}

\subsection{Files Modified}

\subsubsection{\texttt{kernel/proc.h}}
\begin{itemize}
    \item Added 3 MLFQ fields to \texttt{struct proc} (lines 109-112)
\end{itemize}

\subsubsection{\texttt{kernel/proc.c}}
\begin{itemize}
    \item MLFQ configuration and data structures (lines 9-25)
    \item Modified \texttt{scheduler()} function (lines 475-554)
    \item Modified \texttt{yield()} function with demotion logic (lines 585-609)
    \item Modified \texttt{allocproc()} to initialize MLFQ fields (lines 146-149)
    \item Modified \texttt{userinit()} to enqueue init process (lines 258-266)
    \item Modified \texttt{kfork()} to enqueue new processes (lines 339-347)
    \item Modified \texttt{kexit()} to remove process from queue (lines 403-406)
    \item Modified \texttt{sleep()} to remove process from queue (lines 660-663)
    \item Modified \texttt{wakeup()} to re-enqueue processes (lines 691-693)
    \item Implemented \texttt{mlfq\_init()} (lines 868-878)
    \item Implemented \texttt{mlfq\_enqueue()} (lines 880-899)
    \item Implemented \texttt{mlfq\_dequeue()} (lines 901-917)
    \item Implemented \texttt{mlfq\_remove()} (lines 919-945)
    \item Implemented \texttt{mlfq\_tick()} (lines 804-825)
    \item Implemented \texttt{mlfq\_boost()} (lines 827-866)
\end{itemize}

\subsubsection{\texttt{kernel/defs.h}}
\begin{itemize}
    \item Added function declarations (lines 104-109)
\end{itemize}

\subsubsection{\texttt{kernel/main.c}}
\begin{itemize}
    \item Added \texttt{mlfq\_init()} call (line 31)
\end{itemize}

\subsubsection{\texttt{kernel/trap.c}}
\begin{itemize}
    \item Integrated \texttt{mlfq\_tick()} in \texttt{usertrap()} (line 85)
    \item Integrated \texttt{mlfq\_tick()} in \texttt{kerneltrap()} (line 160)
\end{itemize}

\subsubsection{\texttt{Makefile}}
\begin{itemize}
    \item Added \texttt{\_mlfqtest} to UPROGS (line 146)
\end{itemize}

\subsection{New Files}

\subsubsection{\texttt{user/mlfqtest.c}}
\begin{itemize}
    \item Comprehensive test program for MLFQ scheduler
    \item Tests basic scheduling, priority boost, fairness, and responsiveness
\end{itemize}

\subsection{Summary of Changes}

\begin{table}[h]
\centering
\begin{tabular}{@{}ll@{}}
\toprule
\textbf{Component} & \textbf{Changes} \\
\midrule
Data Structures & Added 3 fields to \texttt{struct proc}, 4 queue structures \\
Functions & 6 new MLFQ functions, 2 modified core functions \\
Integration Points & 5 process lifecycle functions modified \\
Initialization & 1 initialization function, 1 call site \\
Testing & 1 test program, test infrastructure \\
\bottomrule
\end{tabular}
\caption{Summary of Code Changes}
\end{table}

% Testing
\section{Testing}

\subsection{Test Program: \texttt{mlfqtest}}

The test program (\texttt{user/mlfqtest.c}) verifies:

\begin{enumerate}
    \item \textbf{Basic Scheduling}: CPU-bound vs I/O-bound process behavior
    \item \textbf{Priority Boost}: Long-running processes get priority boost every 100 ticks
    \item \textbf{Fairness}: Multiple CPU-bound processes get fair CPU time
    \item \textbf{Responsiveness}: I/O-bound processes remain responsive with CPU hogs
\end{enumerate}

\subsection{Test Infrastructure}

Created automated testing tools:
\begin{itemize}
    \item \texttt{test-mlfq.py}: Python script for automated testing
    \item \texttt{run-tests.sh}: Bash script for quick testing
    \item \texttt{generate-report.sh}: Report generation tool
\end{itemize}

\subsection{Running Tests}

\begin{lstlisting}[language=bash]
# Automated testing
./run-tests.sh

# Or manually
make qemu
# In xv6 shell:
mlfqtest
\end{lstlisting}

\subsection{Expected Behavior}

\begin{itemize}
    \item[\textcolor{green!70!black}{$\checkmark$}] I/O-bound processes complete faster (stay at higher priority)
    \item[\textcolor{green!70!black}{$\checkmark$}] CPU-bound processes are demoted (move to lower queues)
    \item[\textcolor{green!70!black}{$\checkmark$}] No process starves (priority boost every 100 ticks)
    \item[\textcolor{green!70!black}{$\checkmark$}] Fair CPU allocation among similar processes
\end{itemize}

% Results and Verification
\section{Results and Verification}

\subsection{Requirements Verification}

All CEP requirements have been satisfied:

\begin{table}[h]
\centering
\begin{tabular}{@{}lll@{}}
\toprule
\textbf{Requirement} & \textbf{Status} & \textbf{Details} \\
\midrule
1. MLFQ Fields in proc.h & \textcolor{green!70!black}{$\checkmark$} & queue\_level, time\_in\_queue, next \\
2. MLFQ Configuration & \textcolor{green!70!black}{$\checkmark$} & 4 queues, time allotments, locks \\
3. mlfq\_init() \& Call & \textcolor{green!70!black}{$\checkmark$} & Implemented, declared, called \\
4. Queue Operations & \textcolor{green!70!black}{$\checkmark$} & enqueue, dequeue, remove \\
5. Scheduler Modification & \textcolor{green!70!black}{$\checkmark$} & Priority-based selection \\
6. Timer Interrupt Handling & \textcolor{green!70!black}{$\checkmark$} & Time tracking, demotion, boost \\
7. Priority Boost & \textcolor{green!70!black}{$\checkmark$} & Every 100 ticks \\
8. Process Lifecycle & \textcolor{green!70!black}{$\checkmark$} & Fork, exit, sleep, wakeup \\
9. Function Declarations & \textcolor{green!70!black}{$\checkmark$} & All in defs.h \\
10. Build System & \textcolor{green!70!black}{$\checkmark$} & mlfqtest in Makefile \\
\bottomrule
\end{tabular}
\caption{CEP Requirements Verification}
\end{table}

\subsection{Build Status}

\textcolor{green!70!black}{\textbf{Kernel compiles successfully}} without errors or warnings

\subsection{Code Quality}

\begin{itemize}
    \item[\textcolor{green!70!black}{$\checkmark$}] Proper locking (no race conditions)
    \item[\textcolor{green!70!black}{$\checkmark$}] Lock ordering to prevent deadlocks
    \item[\textcolor{green!70!black}{$\checkmark$}] Error handling
    \item[\textcolor{green!70!black}{$\checkmark$}] Code comments and documentation
\end{itemize}

\subsection{Test Results}

\begin{itemize}
    \item[\textcolor{green!70!black}{$\checkmark$}] Test program executes successfully
    \item[\textcolor{green!70!black}{$\checkmark$}] All test phases complete
    \item[\textcolor{green!70!black}{$\checkmark$}] No panics or errors observed
    \item[\textcolor{green!70!black}{$\checkmark$}] Expected behavior verified
\end{itemize}

% Conclusion
\section{Conclusion}

\subsection{Achievements}

\begin{enumerate}
    \item \textbf{Successfully implemented MLFQ scheduler} with 4 priority levels
    \item \textbf{Integrated with xv6 process lifecycle} (fork, exit, sleep, wakeup)
    \item \textbf{Implemented priority boost mechanism} to prevent starvation
    \item \textbf{Created comprehensive test suite} for verification
    \item \textbf{All CEP requirements satisfied}
\end{enumerate}

\subsection{Key Features}

\begin{itemize}
    \item \textbf{Priority-based scheduling}: Processes start at highest priority
    \item \textbf{Automatic demotion}: CPU-bound processes move to lower queues
    \item \textbf{Time-based allocation}: Different time allotments per queue
    \item \textbf{Starvation prevention}: Periodic priority boost every 100 ticks
    \item \textbf{I/O optimization}: I/O-bound processes maintain higher priority
\end{itemize}

\subsection{Performance Characteristics}

\begin{itemize}
    \item \textbf{Time Complexity}: O(n) for scheduler (n = number of queues)
    \item \textbf{Space Complexity}: O(1) per process (3 additional fields)
    \item \textbf{Lock Contention}: Minimal (per-queue locks)
\end{itemize}

\subsection{Future Enhancements (Optional)}

\begin{enumerate}
    \item Configurable time allotments and boost interval
    \item Per-CPU queues for multi-core systems
    \item Statistics tracking (queue distribution, wait times)
    \item Debug output for queue states
\end{enumerate}

% Appendix
\section{Appendix}

\subsection{File Locations}

\textbf{Core Implementation}:
\begin{itemize}
    \item \texttt{kernel/proc.h}: Process structure with MLFQ fields
    \item \texttt{kernel/proc.c}: MLFQ implementation (all functions)
    \item \texttt{kernel/defs.h}: Function declarations
    \item \texttt{kernel/main.c}: Initialization call
    \item \texttt{kernel/trap.c}: Timer interrupt integration
\end{itemize}

\textbf{Testing}:
\begin{itemize}
    \item \texttt{user/mlfqtest.c}: Test program
    \item \texttt{test-mlfq.py}: Automated test script
    \item \texttt{run-tests.sh}: Test runner
    \item \texttt{TESTING.md}: Testing documentation
\end{itemize}

\textbf{Documentation}:
\begin{itemize}
    \item \texttt{CEP\_VERIFICATION.md}: Detailed requirement verification
    \item \texttt{MLFQ\_IMPLEMENTATION\_SUMMARY.md}: Implementation overview
    \item \texttt{REQUIREMENTS\_CHECKLIST.md}: Quick checklist
    \item \texttt{CEP\_REPORT.md}: This document
\end{itemize}

\subsection{Key Code Snippets}

\subsubsection{Scheduler Selection Logic}
\begin{lstlisting}
// Iterate through queues from highest to lowest priority
for(int q = 0; q < NMLFQ; q++) {
  acquire(&mlfq[q].lock);
  if(mlfq[q].head != 0) {
    // Find first RUNNABLE process
    // Remove from queue and schedule
  }
}
\end{lstlisting}

\subsubsection{Demotion Logic}
\begin{lstlisting}
if(p->time_in_queue >= time_allotment[p->queue_level]) {
  if(p->queue_level < NMLFQ - 1) {
    p->queue_level++;  // Demote to lower queue
  }
  p->time_in_queue = 0;
}
\end{lstlisting}

\subsubsection{Priority Boost}
\begin{lstlisting}
if(time_since_boost >= BOOST_INTERVAL) {
  time_since_boost = 0;
  mlfq_boost();  // Move all to queue 0
}
\end{lstlisting}

\subsection{Testing Output Example}

\begin{lstlisting}[language=bash, basicstyle=\ttfamily\small]
========================================
   MLFQ Scheduler Verification Tests
========================================

=== Test 1: Basic MLFQ Scheduling ===
Creating processes with different characteristics...
[Process execution with C/I/M indicators]

=== Test 2: Priority Boost Test ===
[Long-running process with periodic boost]

=== Test 3: Fairness Test ===
[Multiple CPU-bound processes]

=== Test 4: Responsiveness Test ===
[I/O-bound process with CPU hog]

========================================
     All MLFQ Tests Completed!
========================================
\end{lstlisting}

\subsection{References}

\begin{itemize}
    \item xv6-riscv source code
    \item Operating Systems: Three Easy Pieces (MLFQ chapter)
    \item Original MLFQ paper by Corbato et al.
\end{itemize}

% Submission Checklist
\section*{Submission Checklist}

\begin{itemize}
    \item[\textcolor{green!70!black}{$\checkmark$}] All code changes implemented
    \item[\textcolor{green!70!black}{$\checkmark$}] All requirements satisfied
    \item[\textcolor{green!70!black}{$\checkmark$}] Code compiles without errors
    \item[\textcolor{green!70!black}{$\checkmark$}] Test program created and working
    \item[\textcolor{green!70!black}{$\checkmark$}] Documentation complete
    \item[\textcolor{green!70!black}{$\checkmark$}] Verification report generated
\end{itemize}

\end{document}

